%% toc: P2 21 21.8
\section{Proofs for B\"uchi's prefix form of the refined reduction formula}
\label{sec:annex-p02-w02c-p02-025-buchi-form}

Programs, configurations and admissibility are those of
Definition~\ref{def:p02-w02c-p02-025-program}; the computation of a program
\(P\) is written \(c_t=(q_t,B_t,h_t)\) for \(t<\ell\), where
\(\ell\in\mathbb N\cup\{\infty\}\) is its length, and \(d_t\) is the move of
step \(t\).  For a formula \(\varphi(x,u,y)\) and elements \(v,v',w\),
\(\mathfrak A\entails\varphi[v,v',w]\) means that \(\varphi\) holds when
\(x,u,y\) take the values \(v,v',w\).  ``Conjunct (\(\beta\)\(n\))'' refers to
the list in Definition~\ref{def:p02-w02c-p02-025-formula}.

\begin{annexitem}{annex:p02-w02c-p02-025-admissible}{thm:p02-w02c-p02-025-admissible}{Elementary bounds for admissible programs}
\begin{proof}
Each step changes the head position by at most one and \(h_0=0\), so
\(h_t\leq t\) for every \(t<\ell\).

For \(t\geq1\) we show \(h_t<t\) by induction.  If \(c_1\) exists, then
\(d_0\neq R\), since (N2) would require \(h_0\leq-2\); and \(d_0\neq L\), since
\(h_0=0\) and the step from \(c_0\) to \(c_1\) is a successor step, which never
moves left from cell \(0\).  So \(h_1=0<1\).  If \(c_{t+1}\) exists and
\(h_t<t\) with \(t\geq1\): when \(d_t=R\), (N2) gives \(h_t\leq t-2\) and so
\(h_{t+1}=h_t+1\leq t-1<t+1\); otherwise \(h_{t+1}\leq h_t<t+1\).

For the tape, \(B_0=\bar0\), and \(B_{t'+1}\) differs from \(B_{t'}\) at most
in cell \(h_{t'}\).  Hence if \(B_t(c)\neq0\), then \(c=h_{t'}\) for some
\(t'<t\), and \(c=h_{t'}\leq t'<t\).  So \(B_t(c)=0\) whenever \(c\geq t\).
\end{proof}
\end{annexitem}

\begin{annexitem}{annex:p02-w02c-p02-025-universality}{thm:p02-w02c-p02-025-universality}{Binary admissible programs simulate Turing machines}
\begin{proof}
Fix a valid code \(e\) and write \(M=M_e=(Q,\mathcal A_{\mathrm{in}},
\mathcal A,\delta,q_0,q_{\mathrm{acc}},q_{\mathrm{rej}})\), with \(m\),
\(\gamma\), \(k=2m+1\geq3\), \(\operatorname{blk}\), \(\operatorname{enc}\) and
\(\operatorname{pos}\) as in Definition~\ref{def:p02-w02c-p02-025-simulation}.
Write the computation of \(M\) on the empty input as
\(C_0,C_1,\ldots\) with \(C_t=(q_t^M,T_t,\eta_t)\), so \(C_0=(q_0,T_\varepsilon,0)\)
with \(T_\varepsilon\) everywhere blank.  For a configuration \(C=(q,T,\eta)\)
of \(M\) with \(q\in Q^{\circ}\) and \((j,\sigma)=\operatorname{pos}(\eta)\), let
\[
 R(C)\defeq\bigl(\mathrm{rd}(q,\sigma,\varepsilon),\ \operatorname{enc}(T),\
 kj\bigr),
\]
a configuration of \(\operatorname{Sim}(M)\).  Two facts about blocks are used.
First, \(f_j=1\) exactly when \(j=0\).  Second, the cells \(j\) and \(-j-1\) of
\(M\) lie in block \(j\) and in no other block, so if \(T'\) differs from \(T\)
only in cell \(\eta\) and \(\operatorname{pos}(\eta)=(j,\sigma)\), then
\(\operatorname{blk}_{T'}(j'')=\operatorname{blk}_T(j'')\) for every
\(j''\neq j\).

\emph{Claim A (start).}  The first step of \(\operatorname{Sim}(M)\) is taken
in state \(\iota\) on the all-zero tape at cell \(0\).  If \(q_0=q_{\mathrm{rej}}\)
it enters \(\bot_1\), and if \(q_0=q_{\mathrm{acc}}\) it enters \(\bot_2\); in
both cases \(M\) halts at once in the same way.  Otherwise it writes \(1\) in
cell \(0\), stays, and enters \(\mathrm{rd}(q_0,+,\varepsilon)\).  The tape is then
\(\operatorname{enc}(T_\varepsilon)\), because \(\gamma(\sqcup)=0^m\) makes
\(\operatorname{blk}_{T_\varepsilon}(0)=1\,0^m\,0^m\) and
\(\operatorname{blk}_{T_\varepsilon}(j)=0^k\) for \(j\geq1\).  Since
\(\operatorname{pos}(0)=(0,+)\), the configuration after one step is \(R(C_0)\).

\emph{Claim B (one step of \(M\)).}  Let \(C=(q,T,\eta)\) with \(q\in Q^{\circ}\),
let \(C'=(q',T',\eta')\) be its successor in \(M\), and let
\((j,\sigma)=\operatorname{pos}(\eta)\) and
\(v=\operatorname{blk}_T(j)=f\,x\,y\).  Started in \(R(C)\), the program
\(\operatorname{Sim}(M)\) passes only through the states
\(\mathrm{rd}\), \(\mathrm{wr}\) and \(\mathrm{tr}\), never moves left from
cell \(0\), and: if \(q'\in\{q_{\mathrm{rej}},q_{\mathrm{acc}}\}\), it enters
\(\bot_1\) or \(\bot_2\) respectively after exactly \(k\) steps; if
\(q'\in Q^{\circ}\), it reaches \(R(C')\) after \(n\) steps with
\(2k-1\leq n\leq 3k-2\).

\emph{Reading.}  For \(0\leq i\leq k-1\), after \(i\) steps the configuration is
\[
 \bigl(\mathrm{rd}(q,\sigma,v_0\cdots v_{i-1}),\ \operatorname{enc}(T),\
 kj+i\bigr).
\]
This
holds for \(i=0\); and for \(i<k-1\) the state has a word of length
\(i<k-1\), the scanned letter is \(\operatorname{enc}(T)(kj+i)=v_i\), and the
transition rewrites \(v_i\) and moves right.

\emph{The decision.}  Step \(k\) is taken at cell \(kj+k-1\) in state
\(\mathrm{rd}(q,\sigma,v_0\cdots v_{k-2})\) and scans \(v_{k-1}\), so the word
formed is \(v\).  It is proper, since \(x=\gamma(T(j))\) and
\(y=\gamma(T(-j-1))\).  If \(\sigma=+\), then \(\eta=j\) and
\(a^{*}=\gamma^{-1}(x)=T(\eta)\); if \(\sigma=-\), then \(\eta=-j-1\) and
\(a^{*}=\gamma^{-1}(y)=T(\eta)\).  So \(\delta(q,a^{*})=(q',b,d)\) is the step of
\(M\) from \(C\): \(T'=T[\eta\mapsto b]\) and \(\eta'=\eta-1\) or \(\eta+1\) for
\(d=L\) or \(R\).  If \(q'\) is \(q_{\mathrm{rej}}\) or \(q_{\mathrm{acc}}\), step
\(k\) enters \(\bot_1\) or \(\bot_2\), and the claim holds.  Otherwise
\(v'=\operatorname{blk}_{T'}(j)\): for \(\sigma=+\),
\(\operatorname{blk}_{T'}(j)=f_j\,\gamma(b)\,\gamma(T(-j-1))=f\,\gamma(b)\,y\), since
\(-j-1\neq j=\eta\); for \(\sigma=-\),
\(\operatorname{blk}_{T'}(j)=f\,\gamma(T(j))\,\gamma(b)=f\,x\,\gamma(b)\).  Step \(k\)
writes \(v'_{k-1}\) in cell \(kj+k-1\) and moves left to
\(kj+k-2\geq kj\), entering \(\mathrm{wr}(v',k-2,q',\sigma',\Delta)\).

\emph{Writing.}  For \(0\leq i\leq k-2\), after \(k+(k-2-i)\) steps the
configuration is \((\mathrm{wr}(v',i,q',\sigma',\Delta),B^{(i)},kj+i)\), where
\(B^{(i)}\) is \(\operatorname{enc}(T)\) with the cells
\(kj+i+1,\ldots,kj+k-1\) holding \(v'_{i+1},\ldots,v'_{k-1}\).  This holds for
\(i=k-2\) by the decision step, and the transition of
\(\mathrm{wr}(v',i,\ldots)\) with \(i\geq1\) writes \(v'_i\) and moves left from
\(kj+i\geq kj+1\).  Step \(2k-1\) is taken in \(\mathrm{wr}(v',0,\ldots)\) at cell
\(kj\) and writes \(v'_0\); afterwards block \(j\) holds \(v'\) and every other
block is unchanged, so by the second fact about blocks the tape is
\(\operatorname{enc}(T')\).

\emph{Moving.}  The table of
Definition~\ref{def:p02-w02c-p02-025-simulation} gives
\(\operatorname{pos}(\eta')=(j+\Delta,\sigma')\), using \(f=f_j\):
\begin{itemize}
\item \(\sigma=+\), \(\eta=j\geq0\).  If \(d=R\), \(\eta'=j+1\geq0\) and
  \(\operatorname{pos}(\eta')=(j+1,+)\).  If \(d=L\) and \(f=0\), then
  \(j\geq1\), \(\eta'=j-1\geq0\) and \(\operatorname{pos}(\eta')=(j-1,+)\).  If
  \(d=L\) and \(f=1\), then \(j=0\), \(\eta'=-1\) and
  \(\operatorname{pos}(\eta')=(0,-)\).
\item \(\sigma=-\), \(\eta=-j-1\).  If \(d=L\), \(\eta'=-(j+1)-1\) and
  \(\operatorname{pos}(\eta')=(j+1,-)\).  If \(d=R\) and \(f=0\), then
  \(j\geq1\), \(\eta'=-j=-(j-1)-1\) and \(\operatorname{pos}(\eta')=(j-1,-)\).  If
  \(d=R\) and \(f=1\), then \(j=0\), \(\eta'=0\) and
  \(\operatorname{pos}(\eta')=(0,+)\).
\end{itemize}
If \(\Delta=0\), step \(2k-1\) stays at \(kj\) and enters
\(\mathrm{rd}(q',\sigma',\varepsilon)\), which is \(R(C')\) after \(2k-1\) steps.  If
\(\Delta=+1\), step \(2k-1\) moves right to \(kj+1\) and enters
\(\mathrm{tr}(R,q',\sigma',k-1)\); the next \(k-1\) steps rewrite the scanned
letters, move right, count \(n\) down to \(1\) and then enter
\(\mathrm{rd}(q',\sigma',\varepsilon)\) at cell \(kj+1+(k-1)=k(j+1)\), which is
\(R(C')\) after \(3k-2\) steps.  If \(\Delta=-1\), then \(f=0\), so \(j\geq1\);
step \(2k-1\) moves left to \(kj-1\geq k-1\), and the next \(k-1\) steps move
left through \(kj-2,\ldots,kj-k=k(j-1)\geq0\) and enter
\(\mathrm{rd}(q',\sigma',\varepsilon)\), which is \(R(C')\) after \(3k-2\) steps.
Every left move in the reading, writing and moving phases starts from a cell
\(\geq1\), which proves Claim~B.

\emph{Claim C (\(\operatorname{Sim}(M)\)).}  If \(q_0\) is halting, Claim~A
settles the matter.  Otherwise, by Claims A and~B and induction on \(t\): for
every \(t\) such that \(C_t\) exists and \(q^M_t\in Q^{\circ}\), the computation of
\(\operatorname{Sim}(M)\) reaches \(R(C_t)\) at a time \(n_t\), with \(n_0=1\)
and \(n_{t+1}\geq n_t+2k-1\), and all its states before \(n_t\) are non-stop;
and if \(C_{t+1}\) exists with \(q^M_{t+1}=q_{\mathrm{rej}}\) (or
\(q_{\mathrm{acc}}\)), it enters \(\bot_1\) (or \(\bot_2\)) at time
\(n_t+k\), its earlier states being non-stop.  Hence if \(M\) rejects,
\(\operatorname{Sim}(M)\) halts in \(\bot_1\); if \(M\) accepts, it halts in
\(\bot_2\); and if \(M\) diverges, the times \(n_t\) are unbounded and no stop
state is ever entered, so the computation of \(\operatorname{Sim}(M)\) is
infinite.  The three outcomes for \(M\) are exhaustive, and the three outcomes
for \(\operatorname{Sim}(M)\) are mutually exclusive, so each implication is an
equivalence.  In each case no step moves left from cell \(0\), so
\(\operatorname{Sim}(M)\) satisfies (N1); the state \(\lambda\) is never
entered.

\emph{Claim D (\(\operatorname{Slow}\)).}  Let \(P'=\operatorname{Sim}(M)\) with
computation \((p_n,B'_n,h'_n)_{n<\ell'}\), so \(p_0=\iota\), a non-stop state.
The computation of \(\operatorname{Slow}(P')\) is
\((a_0,\bar0,0)\), \((a_1,\bar0,0)\), and then, for every \(n<\ell'\) with
\(p_n\) non-stop,
\(c''_{2+2n}=((p_n,0),B'_n,h'_n)\), and \(c''_{3+2n}\) is
\(((p_{n+1},1),B'_{n+1},h'_{n+1})\) if \(p_{n+1}\) is non-stop and
\((p_{n+1},B'_{n+1},h'_{n+1})\) if \(p_{n+1}\) is a stop state.
Indeed the first two steps rewrite the scanned letter and stay; the state
\((p_n,0)\) performs the step of \(P'\) from \((p_n,B'_n,h'_n)\) and records its
target; and \((p_{n+1},1)\) rewrites the scanned letter, stays, and enters
\((p_{n+1},0)\).  So \(\operatorname{Slow}(P')\) halts in \(\bot_1\), halts in
\(\bot_2\), or runs forever exactly as \(P'\) does.  Its moves are those of
\(P'\), made at the times \(2+2n\); all other steps stay.  Hence it never moves
left from cell \(0\), which is (N1).  For (N2), let \(d''_t=R\).  Then
\(t=2+2n\) for some \(n\), and the head position at time \(t\) is \(h'_n\leq n\),
since \(P'\) starts at cell \(0\) and moves at most one cell per step; so
\(h''_t\leq n=(t-2)/2\leq t-2\).

Renumbering the states by a bijection that sends \(a_0,\bot_1,\bot_2\) to
\(0,1,2\) changes no computation, so \(P_e\) is admissible and halts in state
\(1\), halts in state \(2\) or runs forever according as \(M_e\) rejects, accepts
or diverges on the empty input.  These are the conditions \(e\in H_1\),
\(e\in H_2\), \(e\in H_\infty\) of Definition~\ref{def:p02-w02-p02-036-halting}.

\emph{Computability.}  From \(e\in V\) the finite table of \(M_e\) is computed
by decoding.  From the table, \(m\) is found by search, \(\gamma\) is fixed by
listing \(\mathcal A\) in the order given by the code with \(\sqcup\) first and
assigning the binary numerals of length \(m\) in increasing order, and the
finitely many states and transitions of \(\operatorname{Sim}(M_e)\),
\(\operatorname{Slow}(\operatorname{Sim}(M_e))\) and the renumbering are listed
by the clauses of Definition~\ref{def:p02-w02c-p02-025-simulation}.  So
\(e\mapsto P_e\) is computable on \(V\).
\end{proof}
\end{annexitem}

\begin{annexitem}{annex:p02-w02c-p02-025-shape}{thm:p02-w02c-p02-025-shape}{Shape of \(\rho^{\mathrm B}_P\)}
\begin{proof}
Each of (\(\beta\)1)--(\(\beta\)9) is built from atoms \(Hst\), \(T_0st\),
\(Kst\) with \(s,t\in\{x,u,y\}\) and atoms \(Zs\), \(I_ps\) by the connectives
alone; \(\mathrm{Act}\,x\) is a finite disjunction of such atoms, and the
conjunction in (\(\beta\)4) is finite because \(\delta_P\) is a finite table.  No
conjunct contains a quantifier or the equality sign.  So
\(\rho^{\mathrm B}_P=\exists x\,Zx\land\forall x\,\exists u\,\forall y\,\beta_P\)
has the stated form, with the binary symbols \(H,T_0,K\) and the monadic
symbols \(Z,I_0,\ldots,I_r\).  Writing \(\beta_P\) down is a finite
transcription of the table of \(\delta_P\), so \(P\mapsto\rho^{\mathrm B}_P\) is
computable; composed with the computable map \(e\mapsto P_e\) of
Theorem~\ref{thm:p02-w02c-p02-025-universality}, so is
\(e\mapsto\rho^{\mathrm B}_e\).
\end{proof}
\end{annexitem}

\begin{annexitem}{annex:p02-w02c-p02-025-remark}{thm:p02-w02c-p02-025-remark}{The remark: every model satisfies \(\exists y\,Hyy\)}
\begin{proof}
Let \(\mathfrak A\entails\rho^{\mathrm B}_P\).  The first conjunct gives an
element \(z\in Z^{\mathfrak A}\).  Take any element \(v\) (the domain is
nonempty); the second conjunct gives an element \(v'\) with
\(\mathfrak A\entails\beta_P[v,v',w]\) for every \(w\), in particular for
\(w=z\).  The conjunct \(Zy\to(I_0y\land Hyy\land T_0yx)\) of (\(\beta\)1) at
\(y=z\) then yields \(\mathfrak A\entails H[z,z]\), so
\(\mathfrak A\entails\exists y\,Hyy\).
\end{proof}
\end{annexitem}

\begin{annexitem}{annex:p02-w02c-p02-025-intended}{thm:p02-w02c-p02-025-intended}{The intended structure is a model}
\begin{proof}
Let \(D\), \(g\) and the relations be as in
Definition~\ref{def:p02-w02c-p02-025-intended}, and in the halting case let \(s\)
be the halting time, so \(D=\{0,\ldots,s\}\) and \(c_t\) exists exactly for
\(t\in D\).  Since state \(0\) is not a stop state, \(s\geq1\).  We use
Lemma~\ref{thm:p02-w02c-p02-025-admissible} and three preliminary facts.

\emph{(P1) The two parts of \(H\).}  The pairs \((t,h_t)\) satisfy
\(h_t\leq t\).  The pairs \((h_{t+1},t)\) with \(d_t\neq N\) satisfy
\(h_{t+1}<t\): if \(d_t=R\), then \(h_{t+1}=h_t+1\leq t-1\) by (N2); if \(d_t=L\),
then \(h_{t+1}=h_t-1<t\).  Hence for \((v,w)\in H\): if \(w\leq v\) then
\(w=h_v\); and if \(v<w\) then \(v=h_{w+1}\), \(w+1\in D\) and \(d_w\neq N\).

\emph{(P2) The three parts of \(T_0\).}  For \(c<t\) in \(D\):
\((t,c)\in T_0\) iff \(B_t(c)=0\); \((c,t)\in T_0\) iff \(c\leq h_t\); and
\((t,t)\in T_0\) always.

\emph{(P3) Where \(g\) is the successor.}  If \(q_t\notin\{1,2\}\), then
\(t+1\in D\) and \(g(t)=t+1\).  In the halting case \(h_t<s\) for every
\(t\in D\) (as \(h_0=0<s\) and \(h_t<t\leq s\) for \(t\geq1\)), and the
second entry \(w\) of a pair \((v,w)\in H\) with \(v<w\) satisfies \(w<s\).  So
\(g(v)=v+1\) whenever \(v\) is a head position or such a second entry.

Now fix \(v,w\in D\) and check each conjunct of \(\beta_P\) under
\(x\mapsto v\), \(u\mapsto g(v)\), \(y\mapsto w\).

(\(\beta\)1) If \(w\in Z\), then \(w=0\); \(q_0=0\) gives \(I_0(0)\), \(h_0=0\)
gives \((0,0)\in H\), and \((0,v)\in T_0\) by (P2): for \(v=0\) it is diagonal,
for \(v>0\) it holds because \(0\leq h_v\).

(\(\beta\)2) \((v,v)\in T_0\) by (P2).

(\(\beta\)3) If \(v\in Z\) and \(v<w\), then \(w\geq1\); \((w,w)\in H\) would
give \(w=h_w\) by (P1), against \(h_w<w\).

(\(\beta\)4) Suppose \((v,w)\in H\), \(v\not<w\), \(I_p(v)\) with
\(p\notin\{1,2\}\), and \([a](v,w)\).  By (P1), \(w=h_v\); \(p=q_v\); and
\(a=B_v(h_v)\): if \(h_v<v\) by (P2), and if \(h_v=v\) then \(v=0\), the pair is
diagonal, \(a=0\), and \(B_0(0)=0\).  So \(\delta_P(p,a)=(p',b,d)\) is the step
from \(c_v\): \(p'=q_{v+1}\), \(b=B_{v+1}(h_v)\), \(d=d_v\), and
\(g(v)=v+1\) by (P3).  Then \(I_{p'}(v+1)\) holds, and \([b](v+1,h_v)\) holds by
(P2) since \(h_v<v+1\).  If \(d=N\), then \(h_{v+1}=h_v\) and
\((v+1,h_v)\in H\).  For the \(d=R\) clause, the premises at \((w,v)\) give in
the same way \(v=h_w\), \(p=q_w\), \(a=B_w(h_w)\) and \(d_w=R\); by (P3)
\(g(v)=h_w+1=h_{w+1}\), and \((h_{w+1},w)\in H\) because \(d_w\neq N\).  For the
\(d=L\) clause, the premises at \((w,g(v))\) give \(g(v)=h_w\), \(p=q_w\),
\(a=B_w(h_w)\) and \(d_w=L\); the successor \(c_{w+1}\) exists, so by (N1)
\(h_w\geq1\) and \(h_{w+1}=h_w-1\).  If \(v+1\in D\), then \(g(v)=v+1\), so
\(v=h_{w+1}\) and \((v,w)\in H\).  If \(v+1\notin D\), then \(v=s=g(v)=h_w\),
against \(h_w<s\) from (P3).

(\(\beta\)5) If \((w,v)\in H\) and \(w<v\), then by (P1) \(w=h_{v+1}\),
\(v+1\in D\) and \(d_v\neq N\); so \(g(v)=v+1\) and \((v+1,w)=(v+1,h_{v+1})\in H\).

(\(\beta\)6) Suppose \(w\leq v\) and \((v,w)\notin H\), so \(w\neq h_v\) by the
definition of \(H\).  If \(v+1\notin D\), then \(g(v)=v\) and both implications
are trivial.  Otherwise \(g(v)=v+1\), the pair \((v+1,w)\) is below the
diagonal, and \((v+1,w)\in T_0\) iff \(B_{v+1}(w)=0\) by (P2).  Also
\((v,w)\in T_0\) iff \(B_v(w)=0\): for \(w<v\) by (P2), and for \(w=v\) because
the pair is diagonal and \(B_v(v)=0\) by
Lemma~\ref{thm:p02-w02c-p02-025-admissible}.  Since \(w\neq h_v\),
\(B_{v+1}(w)=B_v(w)\), so both implications hold.

(\(\beta\)7) First: if \((w,v)\in H\) and \(w\not<v\), then \(v=h_w\); if
\(h_w<w\), then \((h_w,w)\in T_0\) by (P2), and if \(h_w=w\), then \(w=0\) and the
pair is diagonal.  Second: if \((g(v),w)\in T_0\) and \(g(v)<w\), then
\(g(v)\neq v\) (otherwise \(v=s\geq w\)), so \(g(v)=v+1\leq h_w\) by (P2); then
\(v<w\) and \(v\leq h_w\), so \((v,w)\in T_0\).  Third: if moreover
\((w,v)\in H\), then \(v<w\) puts \((w,v)\) below the diagonal, so \(v=h_w\) by
(P1), against \(v+1\leq h_w\).

(\(\beta\)8) No \(q_t\) equals \(1\).

(\(\beta\)9) \((v,v)\notin K\).  If \((w,v)\in H\), then \(v\) is a head position
\(h_w\) or the second entry of an upper pair, so \(g(v)=v+1\) by (P3) and
\(v<g(v)\).  If \(w<v\), then \(w<g(v)\).  If \(\mathrm{Act}(v)\), then
\(q_v\notin\{1,2\}\), so \(g(v)=v+1\) by (P3).

Hence \(\mathfrak A_P\entails\beta_P[v,g(v),w]\) for all \(v,w\), which gives
\(\mathfrak A_P\entails\forall x\,\exists u\,\forall y\,\beta_P\); and
\(0\in Z\) gives \(\exists x\,Zx\).  In the halting case \(\lvert D\rvert=s+1\).
\end{proof}
\end{annexitem}

\begin{annexitem}{annex:p02-w02c-p02-025-simulation}{thm:p02-w02c-p02-025-simulation}{Simulation lemma for \(\rho^{\mathrm B}_P\)}
\begin{proof}
\emph{The chain.}  Put \(\tau_0\defeq z\).  Given \(\tau_n\), the conjunct
\(\forall x\,\exists u\,\forall y\,\beta_P\) supplies an element \(v'\) with
\(\mathfrak A\entails\beta_P[\tau_n,v',w]\) for all \(w\); let \(\tau_{n+1}\) be
such an element.  Then
\[
 (\ast)\qquad \mathfrak A\entails\beta_P[\tau_n,\tau_{n+1},w]
 \quad\text{for every }n\in\mathbb N\text{ and every element }w .
\]
All references to conjuncts below are to instances of \((\ast)\).

\emph{(G) No backward \(K\).}  We show \(\neg K(\tau_b,\tau_a)\) for \(a\leq b\) by
induction on \(b-a\).  For \(a=b\) this is \(\neg Kxx\) in (\(\beta\)9) at
\(n=a\).  For \(a<b\), if \(K(\tau_b,\tau_a)\), the conjunct \(Kyx\to Kyu\) of
(\(\beta\)9) at \(n=a\), \(w=\tau_b\) gives \(K(\tau_b,\tau_{a+1})\), against the
induction hypothesis for the pair \(a+1\leq b\).

\emph{Induction on \(t<\ell\) for items 1--4.}  For \(t=0\): (\(\beta\)1) at
\(n=0\), \(w=z\) gives \(I_0z\) and \(H(z,z)\), and \(q_0=0\), \(h_0=0\); item~2 is
vacuous; the only \(c\leq0\) is \(c=0=h_0\); and \(T_0(z,z)\) holds by (\(\beta\)2)
at \(n=0\), while \(B_0(0)=0\).

Let \(t+1<\ell\) and assume items 1--4 for \(t\).  Then \(q_t\notin\{1,2\}\).
Put \(a\defeq B_t(h_t)\) and \(\delta_P(q_t,a)=(p',b,d)\), so
\(q_{t+1}=p'\), \(B_{t+1}=B_t[h_t\mapsto b]\) and \(h_{t+1}=h_t+d\).

\emph{Item 2 for \(t+1\).}  By item~1, \(\mathrm{Act}\,\tau_t\); the conjunct
\(\mathrm{Act}\,x\to Kxu\) of (\(\beta\)9) at \(n=t\) gives
\(K(\tau_t,\tau_{t+1})\).  For \(a'<t\), item~2 gives \(K(\tau_{a'},\tau_t)\), and
\(Kyx\to Kyu\) at \(n=t\), \(w=\tau_{a'}\) gives \(K(\tau_{a'},\tau_{t+1})\).  If
\(\tau_a=\tau_b\) for some \(a<b\leq t+1\), then \(K(\tau_a,\tau_a)\), against
\(\neg Kxx\); so \(\tau_0,\ldots,\tau_{t+1}\) are pairwise distinct.

\emph{The head premises.}  By items 3, 1 and~4 for \(t\) and by (G) (as
\(h_t\leq t\)),
\[
 (\pi)\qquad H(\tau_t,\tau_{h_t}),\quad \neg K(\tau_t,\tau_{h_t}),\quad
 I_{q_t}\tau_t,\quad [a](\tau_t,\tau_{h_t}).
\]

\emph{Item 1 and the written letter.}  The first clause of (\(\beta\)4) for
\((q_t,a)\) at \(n=t\), \(w=\tau_{h_t}\) has the premises \((\pi)\) and gives
\(I_{p'}\tau_{t+1}\) and \([b](\tau_{t+1},\tau_{h_t})\).

\emph{Item 4 for \(t+1\).}  For \(c=h_t\): \([b](\tau_{t+1},\tau_{h_t})\) says
that \(T_0(\tau_{t+1},\tau_{h_t})\) holds iff \(b=0\), that is, iff
\(B_{t+1}(h_t)=0\).  For \(c\leq t\) with \(c\neq h_t\): \(\neg K(\tau_t,\tau_c)\)
by (G) and \(\neg H(\tau_t,\tau_c)\) by item~3, so (\(\beta\)6) at \(n=t\),
\(w=\tau_c\) gives \(T_0(\tau_{t+1},\tau_c)\) iff \(T_0(\tau_t,\tau_c)\), iff
\(B_t(c)=0\) by item~4, iff \(B_{t+1}(c)=0\).  For \(c=t+1\): (\(\beta\)2) at
\(n=t+1\) gives \(T_0(\tau_{t+1},\tau_{t+1})\), and \(B_{t+1}(t+1)=0\) by
Lemma~\ref{thm:p02-w02c-p02-025-admissible}.

\emph{Item 3 for \(t+1\): the head is present.}  If \(d=N\), the \(d=N\) clause
at \(n=t\), \(w=\tau_{h_t}\) gives \(H(\tau_{t+1},\tau_{h_t})\), and
\(h_{t+1}=h_t\).  If \(d=R\), the \(d=R\) clause at \(n=h_t\), \(w=\tau_t\) has the
premises \((\pi)\) (its atoms are \(Hyx\), \(\neg Kyx\), \(I_py\), \([a]yx\) with
\(x=\tau_{h_t}\), \(y=\tau_t\)) and gives \(H(\tau_{h_t+1},\tau_t)\).  By (N2),
\(h_t+1\leq t-1<t\), so \(K(\tau_{h_t+1},\tau_t)\) by item~2 for \(t\), and
(\(\beta\)5) at \(n=t\), \(w=\tau_{h_t+1}\) gives
\(H(\tau_{t+1},\tau_{h_t+1})\).  If \(d=L\), then \(h_t\geq1\) by (N1), since
\(c_{t+1}\) exists.  The \(d=L\) clause at \(n=h_t-1\), \(w=\tau_t\) has the
premises \((\pi)\) (its atoms are \(Hyu\), \(\neg Kyu\), \(I_py\), \([a]yu\) with
\(u=\tau_{h_t}\), \(y=\tau_t\)) and gives \(H(\tau_{h_t-1},\tau_t)\).  As
\(h_t-1<t\), \(K(\tau_{h_t-1},\tau_t)\) by item~2 for \(t\), and (\(\beta\)5) at
\(n=t\), \(w=\tau_{h_t-1}\) gives \(H(\tau_{t+1},\tau_{h_t-1})\).

\emph{Item 3 for \(t+1\): the head is unique.}  Write \(h'\defeq h_{t+1}\).  The
conjunct (\(\beta\)3) at \(n=0\), \(w=\tau_{t+1}\) has the premises
\(Z\tau_0\) and \(K(\tau_0,\tau_{t+1})\) (item~2 for \(t+1\)) and gives
\(\neg H(\tau_{t+1},\tau_{t+1})\).  So \(h'\neq t+1\), and \(h'\leq t\) since
\(h'\leq t+1\).  Suppose now \(H(\tau_{t+1},\tau_c)\) for some \(c\leq t\) with
\(c\neq h'\), and let \(i<j\) be the numbers \(c\) and \(h'\) in increasing order;
then \(j\leq t\) and both \(H(\tau_{t+1},\tau_i)\) and \(H(\tau_{t+1},\tau_j)\)
hold.  By (G), \(\neg K(\tau_{t+1},\tau_j)\), so the first conjunct of
(\(\beta\)7) at \(n=j\), \(w=\tau_{t+1}\) gives \(T_0(\tau_j,\tau_{t+1})\).  For
\(l=j-1,j-2,\ldots,i+1\) in turn, the second conjunct of (\(\beta\)7) at \(n=l\),
\(w=\tau_{t+1}\), whose premise \(K(\tau_{l+1},\tau_{t+1})\) holds by item~2 for
\(t+1\) because \(l+1\leq j<t+1\), carries \(T_0(\tau_{l+1},\tau_{t+1})\) to
\(T_0(\tau_l,\tau_{t+1})\).  So \(T_0(\tau_{i+1},\tau_{t+1})\).  The third
conjunct of (\(\beta\)7) at \(n=i\), \(w=\tau_{t+1}\) forbids
\(T_0(\tau_{i+1},\tau_{t+1})\land K(\tau_{i+1},\tau_{t+1})\land
H(\tau_{t+1},\tau_i)\); all three hold (the middle one by item~2, as
\(i+1\leq j\leq t\)).  This contradiction shows \(\neg H(\tau_{t+1},\tau_c)\) for
every \(c\leq t+1\) with \(c\neq h'\).
\end{proof}
\end{annexitem}

\begin{annexitem}{annex:p02-w02c-p02-025-reduction}{thm:p02-w02c-p02-025-reduction}{Refined reduction property for \(\rho^{\mathrm B}_P\)}
\begin{proof}
By (N1), the computation of an admissible \(P\) ends, if at all, in a stop
state; so exactly one of the following holds: \(P\) halts in state \(1\), \(P\)
halts in state \(2\), \(P\) never halts.

\emph{Item 1.}  Let \(P\) halt in state \(1\) at time \(s\), and suppose
\(\mathfrak A\entails\rho^{\mathrm B}_P\).  The conjunct \(\exists x\,Zx\) gives
some \(z\in Z^{\mathfrak A}\), and
Lemma~\ref{thm:p02-w02c-p02-025-simulation} gives \(\tau_s\) with
\(\mathfrak A\entails I_1[\tau_s]\).  But (\(\beta\)8) at \(x=\tau_s\) gives
\(\neg I_1\tau_s\).  So \(\rho^{\mathrm B}_P\) is unsatisfiable.  If \(P\) does not
halt in state \(1\), Lemma~\ref{thm:p02-w02c-p02-025-intended} gives a model.

\emph{Item 2.}  If \(P\) halts in state \(2\) at time \(s\), the intended structure
is a model with \(s+1\) elements by
Lemma~\ref{thm:p02-w02c-p02-025-intended}.  Conversely, let \(\mathfrak A\) be a
finite model.  By item~1, \(P\) does not halt in state \(1\).  If \(P\) never
halted, then \(\ell=\infty\), and item~2 of
Lemma~\ref{thm:p02-w02c-p02-025-simulation} would give pairwise distinct
elements \(\tau_0,\ldots,\tau_t\) for every \(t\), so \(\mathfrak A\) would be
infinite.  Hence \(P\) halts, and in state \(2\).

\emph{Item 3.}  By the trichotomy, \(P\) never halts iff it halts neither in
state \(1\) nor in state \(2\), iff by items 1 and~2 \(\rho^{\mathrm B}_P\) is
satisfiable and has no finite model.
\end{proof}
\end{annexitem}

\begin{annexitem}{annex:p02-w02c-p02-025-exercise}{thm:p02-w02c-p02-025-exercise}{The source's Exercise~2.1.31}
\begin{proof}
For \(e\in V\), Lemma~\ref{thm:p02-w02c-p02-025-shape} gives the form of
\(\rho^{\mathrm B}_e=\rho^{\mathrm B}_{P_e}\) and the computability of
\(e\mapsto\rho^{\mathrm B}_e\).  By
Theorem~\ref{thm:p02-w02c-p02-025-universality}, \(P_e\) is admissible, and
\(e\in H_1\), \(e\in H_2\), \(e\in H_\infty\) according as \(P_e\) halts in state
\(1\), halts in state \(2\), or never halts.  Theorem~\ref{thm:p02-w02c-p02-025-reduction}
turns these three conditions into: \(\rho^{\mathrm B}_e\) is unsatisfiable;
\(\rho^{\mathrm B}_e\) has a finite model; \(\rho^{\mathrm B}_e\) is satisfiable
without finite models, that is,
\(\rho^{\mathrm B}_e\in\mathrm{Sat}\setminus\mathrm{Fin\text{-}sat}=
\mathrm{Inf\text{-}axioms}\).  A sentence is contradictory exactly when it is
unsatisfiable.
\end{proof}
\end{annexitem}

\begin{annexitem}{annex:p02-w02c-p02-025-inseparable}{thm:p02-w02c-p02-025-inseparable}{Trakhtenbrot's theorem for the B\"uchi form}
\begin{proof}
Each \(\rho^{\mathrm B}_e\) belongs to \(\mathcal B\) by
Lemma~\ref{thm:p02-w02c-p02-025-shape}, its monadic symbols being \(Z\) and
\(I_0,\ldots,I_r\).  By the symmetry remark after
Definition~\ref{def:p02-w02-p02-036-inseparable} it suffices to treat each pair
in one order.  Suppose a decidable set \(R\) separates
\(\mathrm{Fin\text{-}sat}(\mathcal B)\) from \(\mathrm{Non\text{-}sat}(\mathcal B)\),
and put \(W\defeq\{e\in V:\rho^{\mathrm B}_e\in R\}\).  \(W\) is decidable,
because \(V\) is decidable, \(e\mapsto\rho^{\mathrm B}_e\) is computable on \(V\)
and \(R\) is decidable.  By Theorem~\ref{thm:p02-w02c-p02-025-exercise}, if
\(e\in H_2\) then \(\rho^{\mathrm B}_e\in\mathrm{Fin\text{-}sat}(\mathcal B)
\subseteq R\), so \(e\in W\); and if \(e\in H_1\) then
\(\rho^{\mathrm B}_e\in\mathrm{Non\text{-}sat}(\mathcal B)\), so
\(\rho^{\mathrm B}_e\notin R\) and \(e\notin W\).  So \(W\) separates \(H_2\) from
\(H_1\), which Lemma~\ref{thm:p02-w02-p02-036-halting} excludes.  The same
argument with \(\mathrm{Inf\text{-}axioms}(\mathcal B)\) in place of
\(\mathrm{Non\text{-}sat}(\mathcal B)\) yields a decidable set separating \(H_2\)
from \(H_\infty\); and with \(\mathrm{Non\text{-}sat}(\mathcal B)\) and
\(\mathrm{Inf\text{-}axioms}(\mathcal B)\) as the pair, one separating \(H_1\)
from \(H_\infty\).  Both are excluded by the same lemma.
\end{proof}
\end{annexitem}
